\documentclass[11pt,a4paper]{article}
\usepackage[T1]{fontenc}
\usepackage{lmodern,amsmath,amssymb,amsthm,mathtools,booktabs,microtype}
\usepackage[margin=27mm]{geometry}
\usepackage[colorlinks=true,linkcolor=blue,urlcolor=blue,citecolor=blue]{hyperref}
\newtheorem{theorem}{Theorem}[section]
\newtheorem{lemma}[theorem]{Lemma}
\newtheorem{proposition}[theorem]{Proposition}
\newtheorem{corollary}[theorem]{Corollary}
\newcommand{\Q}{\mathbb Q}\newcommand{\Z}{\mathbb Z}
\DeclareMathOperator{\Tr}{Tr}
\DeclareMathOperator{\ord}{ord}
\setlength{\emergencystretch}{2em}
\title{A defining polynomial in twenty-five variables:\\
combined scalar and comparison conditions}
\author{}\date{10 November 2026}
\begin{document}\maketitle
\begin{abstract}
Under the four arithmetic inputs of the specified elliptic source,
we construct a computable family of integer polynomials in twenty-five
total variables, each of degree less than $2^{85}$. Every sufficiently
large odd value of one fixed coefficient parameter gives a Diophantine
definition of the integers over the rationals. The construction places
scalar divisibility and elliptic comparisons in two common pole
predicates. The comparison functions enter directly, at two different
local precisions. A linear rational normalization controls the scalar
functions. Its completeness follows by making the common elliptic
index divisible by the selected primes, in addition to the prescribed
reduction orders. A contact-height estimate with exponent $7/3$
then supplies the required global bounds. The aggregate trace-height
theorem has an ineffective additive constant, so the sufficient
coefficient parameter is not numerically certified. Exact source
specifications and independent uniform degree checks accompany the
existence theorem.
\end{abstract}

\section{Statement and arithmetic inputs}

All variable counts include the distinguished rational input. An integer
polynomial $F(X,\mathbf Y)$ defines $\Z$ over $\Q$ if
\[
 x\in\Z\quad\Longleftrightarrow\quad
 \exists\mathbf y\in\Q^{\dim\mathbf Y}\quad F(x,\mathbf y)=0
 \qquad(x\in\Q).
\]
The integer $L$ in the following theorem is a fixed coefficient
parameter, not a quantified variable.

\begin{theorem}\label{thm:main}
Assume the four arithmetic inputs in Section~\ref{sec:inputs}.
There is a computable family of finite polynomial specifications
\[
 F_L\in\Z[X,Y_1,\ldots,Y_{24}],\qquad L\ge3\text{ odd},
\]
with the following properties.
\begin{enumerate}
\item Every member has degree at most
\[
 24072202920766797054476288<2^{85}.
\]
Its number of nonzero monomials is less than $2^{1984}$.
\item There is a fixed $L_0$ such that every odd $L\ge L_0$
defines $\Z$ over $\Q$.
\item The proof does not give an effective sufficient value of $L_0$
or a coefficient-height bound for a member certified to define $\Z$.
\end{enumerate}
No additional $abc$ hypothesis is used.
\end{theorem}

The specification uses finite integer polynomial operations,
determinants, and finite products with explicit index ranges. Its
matrix entries and finite-algebra generators do not count as
quantified variables.

\subsection{The inherited arithmetic inputs}\label{sec:inputs}

Put
\[
 K=\Q(\sqrt2),\qquad d=3-\sqrt2,\qquad D=d^2=11-6\sqrt2,
 \qquad E:y^2=x^3-Dx,
\]
and
\[
 P=\left(\frac{57-28\sqrt2}{8},
          \frac{-660+417\sqrt2}{32}\right).
\]
We use the following inputs, with the scope established in
\cite{Dependencies}.
\begin{enumerate}
\item $2E(K)=\Z P$, and $h(x([n]P))\ge n^2-6$ for every
nonzero integer $n$.
\item The full normalized modular realization for the five contacts is
specified in~\cite{Evaluation,Multiplicity}. Its contact-height
consequence is as follows. If $q=u/v\in\Q$ is in lowest terms,
$v>0$, and $q\notin\{0,1,2401,-343/32\}$, let $M_5(q)$ be the
product of primes $p>76176$ at which at least one of
\[
 v_p(v),\quad v_p(u),\quad v_p(u-v),\quad
 v_p(u-2401v),\quad v_p(32u+343v)
\]
is positive and odd. Then
\begin{equation}\label{eq:contact-height}
                  h(q)<H M_5(q)^{7/3},\qquad H=2^{221234}.
\end{equation}
The normalized modular realization is an explicit hypothesis here.
The exponent $7/3$ is a proved consequence of that realization;
it does not require an additional $abc$ hypothesis.
\item With $B=17b$, the pole predicate consists of
\begin{equation}\label{eq:pole}
 p_0^2=B(1-e)(1+3e),\qquad
 Bq_0^2=(s^2-Be(B-b^2))(s^2+3Be(B-b^2)).
\end{equation}
A rational solution excludes negative odd valuations of $b$ outside
$\{2,3,17\}$. A solution exists whenever the reduced denominator of
$b$ is a square and
\begin{equation}\label{eq:pole-preparation}
 0<b<17,\quad b\in\Q_2^{\times2}\cap\Q_3^{\times2}
                  \cap\Q_{17}^{\times2},\quad
 v_2(b)\ge3,\quad v_3(b)\ge1,\quad v_{17}(b)\ge2.
\end{equation}
We retain the restricted-converse, descent, and simultaneous-realization
inputs supporting these assertions.
\item The simultaneous prime-pattern input supplies the prescribed
prime values and Gaussian generators, their conjugate reduction-order
conditions, arbitrary fixed denominator divisibility, and arbitrary
fixed odd avoidance. It applies to any fixed finite list of rational
contact coefficients after the root points have been chosen.
\end{enumerate}
The last input is the consequence of Sun's Theorem~4.1 stated and proved
in the arithmetic dependency paper. In particular, its denominator
divisibility and avoidance requirements may be enlarged by any fixed
finite set of primes before it is applied.

The aggregate trace-height theorem~\cite{Aggregate}, obtained from
von Buhren's theorem~\cite{VonBuhren}, supplies fixed constants
$C>0$ and $J\ge0$ such that
\begin{equation}\label{eq:trace-height}
 \max_i|n_i|^2\le C h\left(\Tr_{K/\Q}
              \sum_{i=1}^4x([n_i]P)\right)+J
 \qquad(n_i\in\Z\setminus\{0\}).
\end{equation}
The additive constant is ineffective. The companion proof verifies
that the graph of this fixed trace function has no constant abelian
translate through the relevant tuples; this verification includes
translates over the algebraic closure. Thus the constants in
\eqref{eq:trace-height} are uniform in all four labels.

\section{The rational source}

All displayed rational expressions are restricted to their domains.
We require every denominator to be nonzero, every contact value and
pole input to be nonzero, and explicitly require $\lambda\ne0$.
Clearing a rational expression always retains its denominator guard.

\subsection{Nine points and their finite charts}

The nine points are
\[
 A,S,T_1,Z,H_0,U_A,U_S,U_T,U_Z\in E(K).
\]
For each point retain two rational coordinates $h_0,h_1$, and use
\begin{equation}\label{eq:point}
 h=h_0+h_1\sqrt2,\qquad x=h^2,\qquad j^2=h^4-D,
 \qquad y=hj.
\end{equation}
These are exactly the nonzero multiples of $P$. Here is the descent
argument. The point $R_0=(1+\sqrt2,2\sqrt2)$ satisfies $2R_0=P$.
The group input writes every point of $E(K)$ as $mR_0+T$ with
$T\in E[2](K)$. Under the two-isogeny descent map,
$\delta(R_0)=1+\sqrt2$ has opposite signs at the two real embeddings,
whereas the torsion classes $1,-D,d,-d$ have equal signs. A square
abscissa therefore forces $m$ even. Among these torsion classes only
$1$ is a square: the negative classes are excluded by signs and
$N_{K/\Q}(d)=7$ excludes $d$. Conversely every double has square
abscissa. Neither $h=0$ nor $j=0$ occurs in~\eqref{eq:point};
see~\cite{SquareCharts} for the descent-map proof.

For completeness we specify the single finite chart for $j$. Write
\[
 h^4-D=A_0+B_0\sqrt2,
\]
where
\[
 A_0=h_0^4+12h_0^2h_1^2+4h_1^4-11,\qquad
 B_0=4h_0h_1(h_0^2+2h_1^2)+6.
\]
Set $c_0=2A_0+B_0$, $d_0=-A_0-B_0$ and
\[
 A'=6A_0^3-9A_0B_0^2-4B_0^3,\qquad
 B'=-4A_0^3+6A_0B_0^2+3B_0^3.
\]
Adjoin one quartic root $Y$ and reconstruct
\begin{equation}\label{eq:quartic-chart}
 Y^4-4A'Y^2+8(B')^2=0,\qquad
 j=\frac{Y/2+(B'/Y)\sqrt2}{c_0+d_0\sqrt2}.
\end{equation}
If $j=u+v\sqrt2$, the rational component of
$(c_0+d_0\sqrt2)j$ is $2(u^3-2v^3)$, which cannot vanish.
The norm of its multiplier is $2A_0^2-B_0^2\ne0$. The remaining
nonvanishing assertions follow from
\[
 A'^2-2B'^2=(2A_0^2-B_0^2)^2(A_0^2-2B_0^2),
\]
and the discriminant
\[
 2^{15}B'^2(2A_0^2-B_0^2)^4(A_0^2-2B_0^2)^2\ne0.
\]
Indeed $h^4=D$ would imply $N(h)^4=49$, and $B'=0$ would give
a rational root of the Eisenstein polynomial $4T^3-6T-3$.
Thus the point quartic is separable at every rational retained-base
specialization, and its automatic nonzero factors may be omitted from
the guard. These statements are proved in~\cite{QuadraticCharts,SquarePredicate}.

Write the integer labels of $(A,S,T_1,Z)$ as $(n,a,t,b)$, and put
\[
 c=|a-1|,\qquad B_{\mathrm{ind}}=\max(|n|,|a|,|t|,|b|).
\]
The affine formula for $S-P$ excludes $S=\pm P$, so $c>0$. Set
$Q_0=P$, $(Q_1,Q_2,Q_3,Q_4)=(A,S,T_1,Z)$, $U_0=H_0$, and
$(U_1,U_2,U_3,U_4)=(U_A,U_S,U_T,U_Z)$. Define
$V_i=U_i-Q_i$ for $0\le i\le4$. These differences use the usual
affine group formulas and their nonzero equal-abscissa guards.

\subsection{Comparison and scalar functions}

Write $f(Q)=1/x(Q)$ and $z(Q)=-x(Q)/y(Q)$. To avoid confusion
with this local parameter, the scalar coefficient below is denoted
$\rho$. The fourth-order formal-logarithm truncation is simply
\[
                           \ell_4(Q)=z(Q).
\]
Put
\[
 L_i=\frac{\ell_4(U_i)}{\ell_4(H_0)},\qquad
 L'_i=\sigma\left(\frac{\ell_4(V_i)}{\ell_4(V_0)}\right)
 \quad(1\le i\le4),
\]
where $\sigma$ is the nonidentity automorphism of $K$. Let
\[
 \mathcal F(u,v,w)=v^2-u(u-2)vw+(u-2)^2w^2-1,
 \qquad \beta=216X^4+450459.
\]
The sixteen comparison functions $C_1,\ldots,C_{16}$, in order,
are
\begin{equation}\label{eq:comparison-list}
 \begin{gathered}
 (f(U_i))_{i=0}^4,\qquad
 (\sigma f(V_i))_{i=0}^4,\\
 L_i-L'_i\quad(1\le i\le4),\qquad
 L_1-\beta,\qquad \mathcal F(L_2,L_3,L_4).
 \end{gathered}
\end{equation}
The first line contains ten functions. Define the required orders
\begin{equation}\label{eq:comparison-orders}
 a_i=1\quad(1\le i\le10),\qquad a_i=3\quad(11\le i\le16).
\end{equation}
Positive valuation of a reciprocal abscissa is sufficient to put its
point in the formal group. Only the six ratio and residual functions
need precision three.

The six scalar functions $S_1,\ldots,S_6\in\Q$ are the two
rational coefficients, in the basis $1,\sqrt2$, of each of
\begin{equation}\label{eq:scalar-list}
 f(H_0),\qquad f(S-P),\qquad z(S-P)/z(H_0),
\end{equation}
in that order. In particular a lower valuation bound on all six
functions holds at every place of $K$ above the given rational prime.

\subsection{The scalar coefficient and two contact functions}

For fixed odd $L\ge3$, define
\begin{equation}\label{eq:scalar-coefficient}
 \begin{gathered}
 \ell=76249,\qquad
 r=\beta(\beta-1)(\beta-2401)(32\beta+343)(2\beta-1),\\
 C_L=17(\ell^{(L+1)/2}+1),\qquad
 \rho=\frac{102^2\ell^L r}{(1+C_Lr)^2}.
 \end{gathered}
\end{equation}
The normalization theorem and its finite residue certificate
\cite{Normalizer} prove
\begin{equation}\label{eq:forced}
 v_\ell(\rho)=
 \begin{cases}L,&v_\ell(X)\ge0,\\L-20v_\ell(X),&v_\ell(X)<0.
 \end{cases}
\end{equation}
This valuation is always positive and odd. For integral $X$ the
reduced denominator of $\rho$ is a square, $\rho$ is a square in
$\Q_2,\Q_3,\Q_{17}$, and its valuations there satisfy the
preparation bounds in~\eqref{eq:pole-preparation}. In fact
$v_2(\rho)=4$. For every real $X$,
\begin{equation}\label{eq:rho-small}
 0<\rho<\frac{36}{\ell r}<1.
\end{equation}

Let
\[
 \Theta=\{0,-1/2,-1,1/2399,-32/407\},
\]
\[
 \Theta_0=(2+\Theta)\cup(4+\Theta),\qquad
 \Theta_1=(6+\Theta)\cup(8+\Theta),\qquad
 \gamma_0=-1,\quad\gamma_1=-2.
\]
The two contact-and-pole sets are disjoint and contain no zero.
Define
\begin{equation}\label{eq:scales}
 T=\Tr_{K/\Q}\sum_{i=1}^4x(Q_i),\qquad
 A_0=1,\qquad A_1=(1+\rho T^2)^2,
 \qquad t_j=A_j\lambda,
\end{equation}
and
\begin{equation}\label{eq:contacts}
 \delta_j=t_j-\gamma_j,\qquad
 R_j=\frac{\prod_{\theta\in\Theta_j}(t_j-\theta)}{\delta_j^{10}}
 \quad(j=0,1).
\end{equation}
All abscissas in~\eqref{eq:scales} are positive at both real
embeddings, so $T>0$ and $A_1>0$.

Let $\mathcal S$ consist of all primes at most $76176$, together
with the prime divisors of the numerators and denominators of the
22 contact-and-pole values and of their nonzero pairwise differences.
Exact rational arithmetic gives
\begin{equation}\label{eq:exceptional}
 \mathcal S\setminus\{p:p\le76176\}
 =\{1875611,3828397,3982747\}.
\end{equation}
In particular $\ell\notin\mathcal S$.

\subsection{Two pole inputs}

For any nonzero rational coefficient $q$ set
\[
                     \Phi_q(v)=\frac{qv}{1+qv}.
\]
For $j=0,1$, put $\mathcal R_j=\rho R_j$ and
\begin{equation}\label{eq:normalized-functions}
                         G_{j,k}=R_j^3\Phi_{R_j}(S_k).
\end{equation}
Define
\begin{equation}\label{eq:W}
 W_j=
 \prod_{i=1}^{16}\left(1+\mathcal R_j
        \left(\frac{C_i}{\delta_j\mathcal R_j^{a_i}}\right)^2\right)
 \prod_{k=1}^{6}\left(1+\mathcal R_j
        \left(\frac{G_{j,k}}{\mathcal R_j^3}\right)^2\right),
\end{equation}
and
\begin{equation}\label{eq:pole-inputs}
 b_j=\frac{4\rho R_j}{\Tr_{K/\Q}(W_j^{-1})^2}.
\end{equation}
Both inputs receive the pole predicate~\eqref{eq:pole}; both
scalar and comparison functions occur in these two inputs.

Each first conic is represented by
\[
 e=\frac{\tau^2-B}{\tau^2+3B},\qquad
 p_0=\frac{-4B\tau}{\tau^2+3B},\qquad B=17b_j,
\]
with retained coordinates $\tau,s$. The second equation is a
single quadratic equation for $q_0$, guarded nonzero. The inherited
realization provides $p_0q_0s\ne0$ and $e\ne1$ under
\eqref{eq:pole-preparation}, as proved in~\cite{PointSource}.
Indeed its two elliptic points have nonzero abscissas of the prescribed
squareclass and are outside two-torsion, so both ordinates are
nonzero. Rescaling gives the nonzero witnesses in~\eqref{eq:pole};
then $\tau=p_0/(e-1)$ gives the displayed chart. This restriction
changes neither the inherited obstruction nor its required converse.

\section{Local extraction at every required prime}

At primes outside $\mathcal S$, valuations on completions of $K$
are normalized by $v(p)=1$ and have integer values. All arguments
below permit zero test functions, whose valuations are $+\infty$.

\begin{lemma}\label{lem:phi}
Let $e=v(q)+v(v_0)$, allowing $v_0=0$. If $e\ne0$, then
\[
                         v(\Phi_q(v_0))=\max(0,e).
\]
In all cases $v(\Phi_q(v_0))\le\max(0,e)$. When $q$ is a
unit, positive valuation of $\Phi_q(v_0)$ is equivalent to positive
valuation of $v_0$, and these valuations are equal.
\end{lemma}
\begin{proof}
For $e\ne0$, the denominator has valuation $\min(0,e)$.
For $e=0$, the numerator is a unit and the denominator is integral,
so the quotient has nonpositive valuation. These statements also
prove the unit assertion. The assertion for a zero numerator follows
by continuity, with valuation $+\infty$.
\end{proof}

\begin{lemma}\label{lem:trace}
Suppose $v_p(\rho R_j)=\alpha$ is odd and $b_j$ has a pole
certificate. Then $\alpha>0$, and at some place above $p$ each
factor of $W_j$ has valuation greater than $-\alpha/2$. Every
scalar numerator $G_{j,k}$ at that place satisfies
\[
                         v(G_{j,k})>\frac94\alpha.
\]
If $v_p(\delta_j)=0$, then at the same place
\begin{equation}\label{eq:raw-extraction}
                         v(C_i)>(a_i-3/4)\alpha
                         \qquad(1\le i\le16).
\end{equation}
\end{lemma}
\begin{proof}
Each nonzero second summand in~\eqref{eq:W} has odd valuation.
Thus every factor has nonpositive valuation. Set
$q=\max_{\mathfrak p\mid p}v_{\mathfrak p}(W_j)\le0$.
Integrality of local traces gives
$v_p(\Tr(W_j^{-1}))\ge-q$. The pole condition and oddness give
\[
                         1\le v_p(b_j)\le\alpha+2q.
\]
Consequently $\alpha>0$ and $q>-\alpha/2$. At a place attaining
$q$, each factor valuation is at least the sum $q$. For a numerator
$V$ with denominator order $a$, this says
\[
 (1-2a)\alpha+2v(V)>-\alpha/2,
 \qquad v(V)>(a-3/4)\alpha.
\]
Use $a=3$ for each scalar numerator and $a=a_i$ for the comparisons.
The argument also covers a zero numerator, whose factor is $1$.
\end{proof}

\begin{lemma}\label{lem:scalar-odd}
If $p\notin\mathcal S$ and $\mu=v_p(\rho)$ is odd, at least
one $R_j$ is a unit. A source solution then has $\mu>0$ and
\[
 v_p(S_k)>\frac94\mu\quad(1\le k\le6),\qquad
 v_p(c)>\frac94\mu.
\]
In particular $v_p(c)\ge3$ and $v_p(c)\ge2\mu+1$.
\end{lemma}
\begin{proof}
The valuation of $\rho T^2$ is odd. Thus either $A_1\equiv1\pmod p$
or $v_p(A_1)<0$. For each balanced contact function an argument of
positive or negative valuation gives a unit: at negative valuation the
numerator and denominator have the same degree, while at positive
valuation every contact and pole constant is a unit. If $\lambda$
is a unit and $R_0$ is not, its residue is a contact or pole of the
first set. Either $A_1\lambda$ has negative valuation or it has
the same residue as $\lambda$. Disjointness of the two reduced
sets gives a unit $R_1$ in either case. This proves the unit
assertion; compare~\cite{LocalTests}.

At a detector with unit $R_j$, Lemma~\ref{lem:trace} forces
$\mu>0$ and $v(G_{j,k})>9\mu/4$. As $R_j$ is a unit,
Lemma~\ref{lem:phi} gives the same lower bound for the rational
function $S_k$. No integrality assertion about the comparison
functions is needed: their factors have nonpositive valuation, which
is precisely the hypothesis used in the trace argument.

The first four scalar components put $H_0$ and $S-P$ in the formal
group at every place over $p$. Write $H_0=[h]P$, where $h$ is a
nonzero integer. The formal logarithm is an isometry and
$[h](S-P)=[a-1]H_0$, so
\[
 v_{\mathfrak p}(z(S-P)/z(H_0))=v_p((a-1)/h).
\]
The last two scalar components bound the left side by $>9\mu/4$.
Since $v_p(h)\ge0$, the asserted bound for $c$ follows. Integer
valuations and $\mu\ge1$ give $v_p(c)\ge2\mu+1\ge3$.
\end{proof}

\begin{lemma}\label{lem:mixed}
Let $p\notin\mathcal S$ and $\nu=v_p(R_j)>0$ be odd. A source
solution has one of the following two properties:
\begin{enumerate}
\item at some place over $p$, $v(C_i)\ge a_i$ for every $i$;
\item $v_p(c)\ge3$.
\end{enumerate}
\end{lemma}
\begin{proof}
The normalizer $\delta_j$ is a unit. Put $\mu=v_p(\rho)$.
If $\mu$ is odd, Lemma~\ref{lem:scalar-odd} proves the second
alternative, or excludes the solution when $\mu<0$.
If $\mu$ is even, $\alpha=\mu+\nu$ is odd. The trace lemma
forces $\alpha>0$ and gives
\[
 v(C_i)>(a_i-3/4)\alpha\ge a_i-3/4.
\]
As valuations are integers, $v(C_i)\ge a_i$. This includes
negative even $\mu$ and arbitrarily deep common contact valuations.
\end{proof}

\subsection{From comparison functions to integer residuals}

The formal-logarithm estimate needed here has a direct short proof.
If $Q$ is in the first formal neighborhood at a prime
$p>76176$ and $r_Q=v(z(Q))\ge1$, then
\begin{equation}\label{eq:log-error}
 v(\log_E(Q)-\ell_4(Q))\ge5r_Q,
 \qquad \ell_4(Q)/\log_E(Q)\in1+p^{4}\mathcal O.
\end{equation}
Indeed the full logarithm is
$z(Q)\sum_{j\ge0}a_j(Df(Q)^2)^j$, with
$a_0=1$, $a_1=-2/5$ and
$a_j=a_{j-1}(4j-5)/(4j+1)$ for $j\ge2$.
One has $v(Df(Q)^2)=4r_Q$, and the omitted terms start at index
$1$. The denominator of $a_j$ divides $(4j+1)!$. Terms with
$4j+1<p$ have integral coefficients; for the remaining terms,
$j\ge19044$ and
$v(a_j(Df^2)^j)\ge4jr_Q-j\ge3jr_Q\ge4r_Q$.
These bounds tend to infinity and prove~\eqref{eq:log-error}.

Suppose the first alternative of Lemma~\ref{lem:mixed} holds.
The ten reciprocal-abscissa functions in~\eqref{eq:comparison-list} put
$H_0,U_i$ in the formal group at the selected place and $V_0,V_i$
in the formal group at its conjugate. Write $U_i=[u_i]P$ and
$H_0=[h]P$. The full logarithm ratios are
\[
 r_i=u_i/h,\qquad r'_i=(u_i-n_i)/(h-1),
 \qquad (n_1,n_2,n_3,n_4)=(n,a,t,b).
\]
At least one of $h,h-1$ is a $p$-unit. Its ratio is integral;
the cross test and~\eqref{eq:log-error} make the other truncated
ratio integral as well. Both relative errors are then absolute
errors. Modulo $p^3$ one has $L_i\equiv r_i\equiv r'_i$.
Multiplication by $h,h-1$ and subtraction give
\begin{equation}\label{eq:index-congruence}
 L_i\equiv n_i\pmod{p^3}.
\end{equation}
This argument is valid at split and inert primes.

Define the integer residual
\[
 \Delta=\mathcal F(a,t,b).
\]
Then $p^3\mid\Delta$. If $\beta=u/v$ is in lowest terms
with $v>0$, the direct target and~\eqref{eq:index-congruence}
make $\beta$ integral at $p$ and give
$p^3\mid u-nv$.

Let $M_j$ be the product of primes outside $\mathcal S$ at which
$R_j$ has positive odd valuation. For either nonzero integer
$N=\Delta$ or $N=u-nv$, Lemma~\ref{lem:mixed} therefore gives
\begin{equation}\label{eq:residual-divisibility}
                         M_j^3\mid Nc\qquad(j=0,1).
\end{equation}
Primes occurring in both arithmetic coefficients are included in this
divisibility statement.

\section{Heights and soundness}

\subsection{The scalar height and fixed-prime lower bound}

Every prime contributing to $M_5(\beta)$ outside $\mathcal S$
has positive odd valuation in $\rho$. To check this, a negative
valuation of $\beta=216X^4+450459$ at such a prime is a multiple
of four and cannot contribute an odd denominator contact. At a
finite contact exactly one factor of $r$ has positive odd valuation;
the denominator $1+C_Lr$ is a unit. The prime $\ell$ contributes
no contact, by the finite residue certificate for~\eqref{eq:forced}.
Lemma~\ref{lem:scalar-odd} thus gives, with fixed $H_1$,
\begin{equation}\label{eq:scalar-height}
 h(\beta)\le H_1c^{7/9}.
\end{equation}
Omitting the three additional primes in~\eqref{eq:exceptional}
only multiplies the height constant by their product to power
$7/3$. At the fixed prime $\ell$, the same lemma and
\eqref{eq:forced} give
\begin{equation}\label{eq:c-large}
 v_\ell(c)\ge2L+1,\qquad
 c\ge\ell^{2L+1},\qquad L\log\ell\le\frac{\log c}{2}.
\end{equation}

The scalar coefficient has the elementary height bound
\begin{equation}\label{eq:rho-height}
                    h(\rho)\le15h(\beta)+2L\log\ell+C_0,
\end{equation}
where $C_0$ is independent of $X,L$. In fact
$h(r)\le5h(\beta)+\log(2\cdot2402\cdot375\cdot3)$, and
\[
 h(\rho)\le\log(102^2)+L\log\ell+3h(r)+2\log C_L+2\log2.
\]
The bound $C_L\le34\ell^{(L+1)/2}$ proves~\eqref{eq:rho-height}.

\subsection{Height from two disjoint groups of contacts}

For $s\in\{2,4,6,8\}$ the map
\[
                       q_s=2+\frac1{t-s}
\]
sends the contacts $s+\Theta$ to
$\infty,0,1,2401,-343/32$. Its inverse has fixed rational
coefficients. Outside $\mathcal S$, the odd contact radical of
$q_s$ consists exactly of the primes at which $t$ meets one of
these five contacts with odd multiplicity. Within each $R_j$ the
two groups are disjoint modulo every such prime. Consequently their
radicals multiply to $M_j$, and applying the contact-height theorem
to each group and multiplying gives a fixed constant $C_1$ with
\begin{equation}\label{eq:balanced-height}
                         h(t_j)\le C_1 M_j^{7/6}.
\end{equation}
For clarity, each separate application gives
$h(t_j)\le C_sM_s^{7/3}$ with a fixed $C_s$; their product gives
the square of~\eqref{eq:balanced-height}. Fixed exceptional primes
are absorbed into $C_s$.

The nonzero slope gives $A_1=t_1/t_0$. Hence
\[
 h(A_1)\le h(t_0)+h(t_1),\qquad
 h(A_1)=2h(1+\rho T^2)\ge4h(T)-2h(\rho)-2\log2.
\]
Combining~\eqref{eq:trace-height}, \eqref{eq:residual-divisibility},
\eqref{eq:rho-height}, and~\eqref{eq:c-large} yields fixed
constants $C_2\ge1$, $J_2\ge0$, independent of $X,L$, such that
\begin{equation}\label{eq:common-height}
 B_{\mathrm{ind}}^2\le C_2|Nc|^{7/18}
             +C_2c^{7/9}+C_2\log c+J_2
\end{equation}
for each nonzero residual $N$ under consideration.

\subsection{Eliminating the residuals}

Write $B=B_{\mathrm{ind}}$ in this subsection. From~\eqref{eq:c-large},
$B\ge c-1>4$ and $c\le B+1\le2B$. If $\Delta\ne0$, then
$|\Delta|\le5B^4$, so~\eqref{eq:common-height} implies
\[
 B^2\le C_2 10^{7/18}B^{35/18}
       +C_2(2B)^{7/9}+C_2\log(2B)+J_2.
\]
Every term on the right grows more slowly than $B^2$. There is
therefore a fixed threshold, independent of $L$, beyond which this
inequality is impossible. Sufficiently large $c$ forces
$\Delta=0$.

The integer Lucas growth lemma of the source~\cite{PointSource}
then gives
\begin{equation}\label{eq:Lucas-growth}
                              \log B>c-4.
\end{equation}
If $\beta\ne n$, let $N=u-nv\ne0$ with $\beta=u/v$ in lowest
terms. Since $\max(|u|,v)=e^{h(\beta)}$,
\[
                         |Nc|\le4e^{h(\beta)}B^2.
\]
For sufficiently large $c$, the last three terms on the right of
\eqref{eq:common-height} are at most $B^2/2$, because $B\ge c-1$.
Putting $\eta=7/18$ gives
\[
 B^{2-2\eta}\le2C_2 4^\eta e^{\eta h(\beta)},
 \qquad
 \log B\le\frac{\log(2C_2 4^\eta)+\eta H_1c^{7/9}}
                         {2-2\eta}.
\]
This contradicts~\eqref{eq:Lucas-growth} for sufficiently large $c$.
Thus $\beta=n\in\Z$.

For rational $X$, integrality of $216X^4+450459$ is equivalent
to $X\in\Z$: in a reduced fraction $X=u/v$ it implies
$v^4\mid216=2^3 3^3$, and hence $v=1$.
All sufficient lower bounds on $c$ above are fixed independently
of $L$. Equation~\eqref{eq:c-large} therefore proves soundness
for every sufficiently large fixed odd $L$. The only ineffective
constant used here is inherited from~\eqref{eq:trace-height}.

\section{Completeness}

Fix $X\in\Z$ and any fixed odd $L\ge3$. We prove completeness
for this value of $L$, without a lower-bound hypothesis on $L$.
The order of choices below is needed because the prime-pattern
theorem is applied only after all roots, scales, local neighborhoods,
and bounds are fixed.

\subsection{Preparing the scalar functions}

Let $\mathcal P$ be the union of the finite prime support of $\rho$,
the exceptional set $\mathcal S$, and the primes $2,3,17$.
For the shared point
$H_0=[h_0]P$, choose a positive $h_0$ divisible by $64$ so that
$H_0$ belongs to sufficiently deep formal subgroups at these finite
places and is close to the identity at both real embeddings.
After imposing the finite divisibility, such real returns are possible
in the compact closure of the cyclic subgroup. They do not require
this closure to be the full product of the two real elliptic groups.

Next choose a positive $c$ divisible by $6$ and by any further
required finite modulus, with $[c]P$ still closer to the identity
at all these places. Set $a=c+1$ and $S=[a]P$. This makes the
six rational components in~\eqref{eq:scalar-list} arbitrarily small
simultaneously in the real absolute value and at the specified finite
primes. The ratio is controlled because $H_0$ has already been
fixed and has nonzero local parameter. The finite-place assertion
is exactly the scalar realization in the source, including ramified
and bad primes; its proof uses sufficiently deep formal subgroups
and continuity~\cite{Glazer}. The real assertion follows from
the same continuity and positive return times in the compact cyclic
closure.

For the real preparation below, choose the scalar components small
enough that all six scalar factors in~\eqref{eq:W} are as close
to $1$ as required whenever $1/2<R_j<2$. At a fixed finite prime,
choose them small enough for the same conclusion whenever $R_j$
is sufficiently close to $1$. These choices are possible: as
$S_k\to0$, each scalar second summand tends to zero, with
$\rho$ fixed. No condition involving the as yet unchosen slope
is needed.

Set $n=\beta$. The explicit Lucas construction takes
\[
 U_0=0,\quad U_1=1,\quad U_{m+1}=aU_m-U_{m-1},\quad
 k=2(a-2),\quad b=U_k/(a-2),\quad t=U_{k+1}.
\]
The congruence $U_m\equiv m\pmod{a-2}$ shows that $b$ is an
integer; the Lucas identity gives $\mathcal F(a,t,b)=0$.
As $a$ is odd and $a\equiv1\pmod3$, the sequence modulo two
has period $0,1,1$ and $k\equiv1\pmod3$, so $t,b$ are positive
and odd. The four root points and both positive square scales
in~\eqref{eq:scales} are now fixed.

Temporarily set $H=[h]P$ and $U_i=[hn_i]P$. Restrict the shared
index to a fixed congruence class and a compact real neighborhood
on which all prepared scalar bounds persist. Choose $|S_k|<1/4$
throughout that neighborhood. Then $1+R_jS_k>1/2$ whenever
$1/2<R_j<2$, so every scalar normalization denominator is nonzero.
The original point, difference, and ratio functions have only finitely
many poles under each fixed multiplication and translation map.
Avoiding those poles and shrinking the neighborhoods bounds every
raw comparison function. For example, equal-abscissa failures of the
chosen differences are absent for $h>2$ and nonzero $n_i$.
The fixed congruence class is taken modulo a multiple of $64$ and
of the order of the component group of the compact closure of
$\Z P$ in the product of the two real elliptic groups. This component
group is finite because a closed subgroup of a compact real Lie
group has finitely many components.
There are no additional comparison normalization denominators.

\subsection{Choosing the slope and comparison primes}

Apply the inherited simultaneous prime-pattern theorem to the rational
coefficients $-\theta/A_j$, for both $j$ and all
$\theta\in\Theta_j$. Repeated coefficients are listed once,
retaining their incidences. In the notation of its rational-contact
form, choose
\[
                    \lambda=u/v,\qquad v=D_0^2w^2,
\]
so that all $u-(\theta/A_j)v$ are prescribed positive primes.
Require these primes to avoid every previously fixed prime, the
numerator and denominator support of the scales, and all contact
differences. Prescribe enough divisibility of $w$ that $\lambda$
has arbitrarily negative valuation at each prime in $\mathcal P$.
The same theorem supplies an arbitrarily large positive real slope.
The rational-contact and avoidance lemmas of~\cite{Dependencies}
prove these assertions from the stated prime-pattern input.

As $\lambda$ tends to infinity at any of the preparation places, both
$R_j$ tend to $1$ and $\delta_j^{-1}$ tends to zero. The scalar
factors are close to $1$ by their preparation. The comparison
factors tend to $1$ because of $\delta_j^{-1}$, uniformly on
the already bounded comparison neighborhoods. Thus all required
finite and real closeness conditions on $W_j$ can be imposed
before the prime-pattern theorem is used.

We impose one further size requirement before choosing the primes:
\begin{equation}\label{eq:prime-cluster}
 \lambda>4\max_{j,\theta\in\Theta_j}(\theta/A_j),\qquad u\ge25.
\end{equation}
All contacts here are positive. Each selected prime
$p_i=u-(\theta/A_j)v$ therefore lies in $(3u/4,u)$. The full
order of the reduced elliptic curve at either conjugate of any selected prime
$p_k$ is even and at most
\[
 p_k+1+2\sqrt{p_k}<u+1+2\sqrt u<3u/2.
\]
Evenness follows from the full rational two-torsion. The prescribed
annihilating moduli divide these group orders. The last inequality holds for $u\ge25$.
If any selected odd prime $p_i$ divided one of these even orders,
that order would be at least $2p_i>3u/2$, a contradiction.
Thus each selected prime is coprime to every prescribed reduction
order. All of them also avoid the previously fixed modulus.

The conjugate reduction orders give a common arithmetic progression
for $h$: the primary orders divide $h$, the opposite orders divide
$h-1$, and $h$ remains in its previously fixed congruence class.
Their odd coprimality and two-adic compatibility are proved
in~\cite{Dependencies}. By the preceding paragraph, we may add
\begin{equation}\label{eq:extra-divisibility}
                         h\equiv0\pmod{\prod_i p_i}.
\end{equation}
Positive indices in the resulting progression meet the prescribed
real neighborhood. Indeed the initial fixed modulus annihilates the component group
of the compact cyclic closure. Its identity component is a compact
real torus, on which multiplication by every positive integer is
surjective. Every compatible refinement therefore has the same
closure inside the prescribed component; see~\cite{Dependencies}.
Choose such an $h>2$, and set $H_0=[h]P$, $U_i=[hn_i]P$.

At a selected comparison prime $p$, the argument in its detector
meets exactly one contact. Thus $v_p(R_j)=1$, and $\rho,\delta_j$
are units. The reduction-order conditions put the relevant $U_i,V_i$
in their first formal neighborhoods. Consequently the first ten
comparison functions have valuation at least $2$. The exact full
ratios are the integers $n_i$, so~\eqref{eq:log-error} gives
valuation at least $4$ for the last six comparisons. In particular
\[
 v\left(\mathcal R_j
       (C_i/(\delta_j\mathcal R_j^{a_i}))^2\right)\ge3
                         \qquad(1\le i\le16)
\]
at the designated place. All comparison factors are units congruent
to $1$ there.

It remains to check the linear scalar denominators at these new
primes. Before prime selection exclude the finite set where the
coordinates of $P$ or $S-P$, or their nonzero local parameters,
are not units. At the primary place, $[h/p]P$ is in the formal
group: division by $p$ preserves annihilation modulo its reduction
order, since that order is coprime to $p$. Then
$H_0=[p]([h/p]P)$ has local-parameter valuation at least two.
At the opposite place $H_0$ reduces to $P$ and its local parameter
is a unit. The prime is split in $K$, as specified by the inherited
prime pattern. Its two rational components of $f(H_0)$ are units,
because the two local values have respectively positive and zero
valuation. The rational components of $f(S-P)$ are integral.
Finally the two local values of $z(S-P)/z(H_0)$ have valuations
at most $-2$ and exactly $0$. Its two rational components therefore
both have valuation at most $-2$. In these assertions $2$ and
$\sqrt2$ are units, so forming half the sum and the scaled
half-difference preserves the smaller of two unequal valuations.

Every scalar component $S_k$ consequently has either nonnegative
valuation or valuation at most $-2$. As $v_p(R_j)=1$, the product
$R_jS_k$ never has valuation zero. Lemma~\ref{lem:phi} shows
that $\Phi_{R_j}(S_k)$ is integral. Hence $v_p(G_{j,k})\ge3$,
and every scalar second summand in~\eqref{eq:W} has valuation
at least $1$. All scalar factors are units congruent to $1$ at
both places. This argument allows arbitrarily deep formal reduction
of $P$ and does not require a bound on elliptic Wieferich valuations.

\subsection{Preparing the pole inputs}

The factorization of the contact numerators gives the exact squareclass
description
\[
 R_j=\left(\frac{A_j}{v\delta_j}\right)^{10}
                  \prod_{\theta\in\Theta_j}
                    \bigl(u-(\theta/A_j)v\bigr).
\]
Thus each $R_j$ is a positive rational square times its ten selected
primes. Their odd valuations are positive and equal to one, and they
are disjoint from the fixed support of $\rho$. At those primes
$W_j^{-1}\equiv1$ at the designated place. At its conjugate,
each factor of $W_j$ has nonpositive valuation: it is either a
unit congruent to $1$ or has negative valuation, since its second
summand has odd valuation. Thus $W_j^{-1}$ is congruent to $1$
or $0$ there. The trace is consequently congruent to $2$ or $1$
modulo $p$ and is a unit, so the odd valuation of $b_j$ is positive.
At primes in the support of
$\rho$, both $R_j$ and $W_j$ have been made arbitrarily close
to $1$. Every odd valuation of $\rho R_j$ is therefore nonnegative
in $b_j$. At all remaining primes its valuation is even and remains
even after division by the square trace. The reduced denominator
of each $b_j$ is a square.

At $2,3,17$, both $R_j$ and $W_j$ can be as close to $1$ as
needed. Each $R_j$ is then a local square, and the trace tends
to $2$. The squareclass of $b_j$ agrees with that of $\rho$.
At two the factor $4$ in~\eqref{eq:pole-inputs} cancels the
valuation of the limiting trace square, giving $v_2(b_j)=4$.
At three and seventeen the required depth is unchanged. Thus
all finite conditions in~\eqref{eq:pole-preparation} hold.

At each real embedding arrange $1/2<R_j<2$ and $1\le W_j<2$.
The latter inequality follows by making the finitely many positive
second summands sufficiently small. The trace is then at least
$1$, so
\[
                         0<b_j<8\rho<17.
\]
The scalar normalization denominators are nonzero by the prepared
bound $|R_jS_k|<1/2$. All products and traces are nonzero by
positivity. The inherited converse supplies the two pole witnesses
with the nonzero chart conditions already described. All nine point
charts and the guarded differences are available, since the chosen
labels are nonzero, $a>1$, and $h>2$. This completes the rational
source construction for every integral input and every fixed odd
$L\ge3$.

\section{A single polynomial and its degree}

The retained rational coordinates are
\begin{center}
\begin{tabular}{lr}
\toprule
Coordinates & Number\\
\midrule
Input & 1\\
Nine point-coordinate pairs & 18\\
Nonzero slope & 1\\
Two pole-parameter pairs & 4\\
\midrule
Total before finite-root elimination & 24\\
\bottomrule
\end{tabular}
\end{center}
There are nine quartic point roots and two quadratic pole roots.
Their finite-algebra rank is $4^9 2^2=2^{20}$. Both pole roots
are guarded nonzero. The separating-element theorem~\cite{Active}
eliminates all eleven finite roots and retains every nonautomatic
guard with one additional rational coordinate. This gives exactly
twenty-five total polynomial variables.

We specify that elimination explicitly. Let $B_0$ be the algebra of
the nine point roots, $A_0$ the algebra after adjoining the two pole
roots $q_1,q_2$, and write $q_j^2=c_j\in B_0$ after rational
square clearing. Put $b_0=2^{18}$ and $n_0=2^{20}$. The source
guard has the form $G=gq_1q_2$, $g\in B_0$, and there are no
additional closed residual equations. Set
\[
 T_0=g^2c_1c_2,\qquad
 \det_{B_0}(I-vM_{T_0})=\sum_{j=0}^{b_0}a_jv^j,\qquad a_0=1.
\]
Let $\theta_1,\ldots,\theta_m$, $m=11$, be the point and pole
generators, and let $\delta_i$ be the discriminant of the $i$th
point quartic. These discriminants never vanish at rational base
specializations. Define
\[
 R=n_0\left(1+\sum_{i=1}^{m}\sum_{u,v=1}^{n_0}
              (M_{\theta_i}^{A_0})_{uv}^2
            +\sum_{u,v=1}^{b_0}(M_g^{B_0})_{uv}^2\right).
\]
For $1\le r\le b_0$, put
\begin{align*}
 c_r&=a_r,&h_r(Z)&=\sum_{j=0}^r a_jZ^{r-j},&E_r&=c_r-h_r(T_0),\\
 D_r&=c_r^2\prod_{i=1}^9\delta_i^2,&s_r&=3r+108,\\
 K_r&=(2R)^{s_r}(1+c_r^2)\prod_{i=1}^9(1+\delta_i^2),&
 C_r&=D_r+4RK_r,\\
 \beta_r&=D_rE_r\left(C_r^m+
             \sum_{i=1}^mD_r^iC_r^{m-i}\theta_i\right),&
 H_r(U)&=\det_{A_0}(I-UM_{\beta_r}).
\end{align*}
The required integer polynomial is
\begin{equation}\label{eq:finite-product}
 F_L=H_{b_0}(U)\prod_{r=1}^{b_0-1}
            \left(H_r(U)^2+\sum_{j=r+1}^{b_0}a_j^2\right).
\end{equation}
The only additional rational coordinate is $U$.

The proof in~\cite{Active} handles every base specialization.
At the actual degree of $\det(I-vM_{T_0})$, the element $E_r/c_r$
is the idempotent selecting exactly the factors where the guard is
nonzero. The coefficient $c_r$ and the nonzero quartic discriminants
give lower bounds for the separations of the finite roots. The
powers of $C_r/D_r$ distinguish the surviving characters, and
multiplication by $E_r$ kills all other factors, including their
nilpotents. A rational root of the corresponding $H_r$ is therefore
equivalent to a guarded rational character of the finite algebra.
The squared tail coefficients exclude incorrect degree cases; when
$c_r=0$ one has $H_r=1$. Thus~\eqref{eq:finite-product} has
exactly the rational projection of the guarded source, including
specializations with an empty nonzero part or complex residue factors.

The certified weights of the nine point generators are $6$; the
two pole-generator weights are $287146,1541546$. The full
guard and point-guard weights are $2231217,402525$.
The separating elements have weight at most $W_1r+W_0$, where
\[
 W_1=334007250,\qquad W_0=8083776298.
\]
The principal-minor degree estimate, summed over every factor, is
\[
 \deg F_L\le n_0\bigl(W_1b_0^2+(W_0+1)(2b_0-1)\bigr)
       =24072202920766797054476288<2^{85}.
\]
The uniform checker treats $\ell^L$ and $\ell^{(L+1)/2}$ as
independent coefficient symbols of weight zero. Consequently this
degree bound holds for every fixed odd $L$. Integer specialization
cannot increase it. Equations~\eqref{eq:c-large} and
\eqref{eq:common-height}, followed by completeness, now prove
Theorem~\ref{thm:main}.

There is also a uniform bound on the number of nonzero monomials.
Let $\mathcal D=24072202920766797054476288$. From
\eqref{eq:finite-product},
\[
 \deg_U F_L\le\mathcal E=n_0(2b_0-1)=549754765312.
\]
Counting monomials of total degree at most $\mathcal D$ in the
25 variables, with the $U$ exponent at most $\mathcal E$, gives
\[
 \#\operatorname{supp}(F_L)
 \le\binom{\mathcal D+25}{25}
      -\binom{\mathcal D-\mathcal E+24}{25}<2^{1984}.
\]
The last inequality is checked by exact integer arithmetic. It is a
support upper bound, not an assertion of the expanded monomial count.

\section{Exact specifications and reproducibility}

The directory \nolinkurl{computation/few_variables/merged25_v19} contains
the source builder, the guarded rational source, its integer polynomial
specification, the independent source checker, and the uniform degree
certificate. Its finite-root elimination is in the subdirectory
\nolinkurl{active25}. The source has $5860$ rational
coefficient nodes, $2239$ integer coefficient nodes, and $136$
atomic guard factors. The checks include the nonzero slope, the ten raw reciprocal-abscissa tests, the six
order-three residual tests, every linear scalar denominator, the two pole charts,
the guard factorization, and the generator dependencies. The active
specification has $262144$ finite-product factors. These are
specification counts; the expanded monomial count is not asserted.

The supplied numerical fixture uses $L=3$. It checks the source
and finite-polynomial construction for a concrete coefficient
specialization; it is not certified to satisfy the ineffective
sufficient threshold. The defining-polynomial conclusion concerns
every sufficiently large odd member of the family. All programs
for these source and degree checks use the Python standard library.

\begin{thebibliography}{99}
\interlinepenalty=10000
\bibitem{Active}
\emph{Separating the Nonzero Part of a Finite Algebra},
companion paper, 9 November 2026,
\nolinkurl{active_finite_algebra_20261109_v15.tex}.
\bibitem{Aggregate}
\emph{Aggregate trace heights of elliptic multiples},
companion paper, 8 November 2026,
\nolinkurl{aggregate_trace_height_20261108_v9.tex}.
\bibitem{Dependencies}
\emph{Arithmetic dependencies of the construction with 29 variables},
companion paper, 8 November 2026,
\nolinkurl{arithmetic_dependencies_20261108_v8.tex}.
\bibitem{Evaluation}
\emph{Evaluation of Differentials and Heights of Modular Jacobians},
companion paper, 9 November 2026,
\nolinkurl{evaluation_of_differentials_20261109_v17.tex}.
\bibitem{Glazer}
E.~Glazer, \emph{Diophantine definition of the integers over the
rationals}, Version~9, \nolinkurl{manuscripts/quartic_v9.tex},
commit \nolinkurl{14e8f601fe0d9173d4a5f084aedc3e07066a7688}.
\url{https://github.com/elliotglazer/diophantine-z-q}.
The inherited restricted-converse correction is identified
in~\cite{Dependencies}; the scalar realization is
\nolinkurl{v5:scalar32-bridge}, and the formal-logarithm argument is
\nolinkurl{lem:v5-unramified-log}.
\bibitem{LocalTests}
\emph{Rational Normalization of Local Valuation Tests},
companion paper, 10 November 2026,
\nolinkurl{rationally_normalized_local_tests_20261110.tex}.
\bibitem{Multiplicity}
\emph{Multiplicity and Contact-Height Estimates},
companion paper, 10 November 2026,
\nolinkurl{multiplicity_and_contact_heights_20261110_v19.tex}.
\bibitem{Normalizer}
\emph{Scalar contact normalization by finite-field avoidance},
companion paper, 8 November 2026,
\nolinkurl{scalar_normalization_20261108_v9.tex}.
\bibitem{PointSource}
\emph{A defining polynomial in twenty-seven variables:
an existence theorem with a uniform degree bound},
companion source paper, 9 November 2026,
\nolinkurl{few_variables27_nontorsion_20261109_v15.tex}.
\bibitem{QuadraticCharts}
\emph{A single quartic chart for square roots in a quadratic extension},
companion paper, 8 November 2026,
\nolinkurl{quadratic_square_chart_20261108_v9.tex}.
\bibitem{SquareCharts}
\emph{Square abscissas and direct charts for the elliptic subgroup},
companion paper, 8 November 2026.
\bibitem{SquarePredicate}
\emph{Optimal polynomial square tests in quadratic number fields},
companion paper, 8 November 2026,
\nolinkurl{quadratic_square_predicates_20261108_v10.tex}.
\bibitem{Sun}
Z.-W.~Sun, simultaneous prime-pattern theorem,
Theorem~4.1 of arXiv:2610.20691v2, as pinned in~\cite{Dependencies}.
\bibitem{VonBuhren}
J.~von Buhren,
\emph{Borne de hauteur semi-effective pour le probl\`eme de
Mordell--Lang dans une vari\'et\'e ab\'elienne},
J.~Th\'eor. Nombres Bordeaux \textbf{29} (2017), no.~1,
307--320, Theorem~1.1.
\url{https://www.numdam.org/item/10.5802/jtnb.981.pdf}.
\end{thebibliography}
\end{document}
